\section{The full-board limit}\label{sec:limite}

The bounds in the preceding section keep the rescaled factors bounded and bounded away from zero. It remains to identify their limit.
We first determine the optimal density of mass one. The row and column normalizations then allow us to recover its individual factors.

\begin{proposicion}[The problem at mass one]\label{prop:masa_uno} The supremum of $s$ over $\cA_1$ is finite and is attained at a
unique density $\nu_1$ of finite entropy. Moreover, $\sopt(\rho)\to\sopt(1)$ as $\rho\uparrow1$. \end{proposicion}

\begin{proof} By Remark~\ref{obs:slater}, $\nu^\sharp\in\cA_1$ and $s[\nu^\sharp]>-\infty$, and by Lemma~\ref{lem:elementales}(c),
$s\le6/e$ on $\cA_1$. The proof of existence and uniqueness is the same as in Step~1 of the proof of Theorem~\ref{thm:max}, which does
not use $\rho<1$. No Euler equation at mass one is needed. For the limit superior, let $\rho_j\uparrow1$. The entropies of the
$\nu_{\rho_j}$ are bounded (since $\sopt(\rho_j)\ge s[\rho_j\nu^\sharp]$, which is bounded below), every subsequence has a further
subsequence that converges weakly in $L^1$ to some $\bar\nu\in\cA_1$ (as in the proof of Proposition~\ref{prop:continuidad}), and weak
upper semicontinuity gives $\limsup\sopt(\rho_j)\le s[\bar\nu]\le\sopt(1)$. For the limit inferior, $\rho\nu_1\in\cA_\rho$ and \[
 \sopt(\rho)\ge s[\rho\nu_1]=-3\rho-\rho\log\rho-\rho\int\nu_1\log\nu_1+\sum_f\int_{I_f}G(\rho\,m_f[\nu_1]),
\] which tends to $s[\nu_1]=\sopt(1)$ by dominated convergence, since $G$ is continuous and bounded. \end{proof}

\begin{proposicion}[Convergence in $L^2$]\label{prop:conv_L2} As $x\downarrow0$, $\nu_{1-x}\to\nu_1$ in $L^2([0,1]^2)$ and
$c_\nu\le\nu_1\le C_\nu$. Moreover, $u_{D,x}\to u_{D,0}=1-m_D[\nu_1]$ and $u_{A,x}\to u_{A,0}=1-m_A[\nu_1]$ in $L^2$, $b_x\to
b_0\in[\frac12,b_+]$, and $r_{R,x}\to r_{R,0}$, $r_{C,x}\to r_{C,0}$ in $L^2$, with $\nu_1=b_0\,r_{R,0}\,r_{C,0}\,u_{D,0}\,u_{A,0}$
almost everywhere. \end{proposicion}

\begin{proof} By Proposition~\ref{prop:cotas}, the $\nu_\rho$ with $\rho\ge\frac12$ are bounded by $C_\nu$, so every sequence
$\rho_j\uparrow1$ has a subsequence converging weakly in $L^2$ to some $\bar\nu\in\cA_1$ with $c_\nu\le\bar\nu\le C_\nu$. By weak upper
semicontinuity and Proposition~\ref{prop:masa_uno}, $s[\bar\nu]\ge\limsup\sopt(\rho_j)=\sopt(1)$, and by uniqueness $\bar\nu=\nu_1$.
Thus $\nu_\rho\rightharpoonup\nu_1$ and $\nu_1$ inherits the bounds. The convergence is strong. Indeed, $-z\log z$ has second
derivative $-1/z\le-1/C_\nu$ on $(0,C_\nu]$ and the other terms of $s$ are concave, so \[
 s\Bigl[\frac{\nu_\rho+\nu_1}2\Bigr]\ge\frac{\sopt(\rho)+\sopt(1)}2+\frac{\|\nu_\rho-\nu_1\|_2^2}{8C_\nu}.
\] The left-hand side is at most $\sopt(\frac{1+\rho}2)\to\sopt(1)$, and hence $\|\nu_\rho-\nu_1\|_2\to0$.

The projections are bounded operators from $L^2([0,1]^2)$ to $L^2(I_f)$ by Cauchy--Schwarz on fibers of length at most one. Hence
$u_{D,x}$ and $u_{A,x}$ converge in $L^2$, and the limits retain the bound $\kappa\le u\le1$. Lifting a line profile to the square does
not increase its $L^2$ norm, and the two-sided bounds ensure that products and quotients preserve convergence in $L^2$. Thus $K_x\to
K_0=u_{D,0}(r-c)\,u_{A,0}(r+c)$ and $V_x=\nu_{1-x}/K_x=b_x\,r_{R,x}\,r_{C,x}$ converges in $L^2$ to $V_0=\nu_1/K_0$. The identities
\begin{equation}\label{eq:identidades_V}
 b_x=\int_{[0,1]^2}V_x,\qquad r_{R,x}=T_RV_x/b_x,\qquad r_{C,x}=T_CV_x/b_x,
\end{equation} which follow from $\int r_R=\int r_C=1$, give $b_x\to b_0=\int V_0\in[\frac12,b_+]$ and the convergence in $L^2$ of
$r_{R,x}$ and $r_{C,x}$ to $r_{R,0}=T_RV_0/b_0$ and $r_{C,0}=T_CV_0/b_0$, respectively. Passing to the limit in
$V_x=b_x\,r_{R,x}\,r_{C,x}$ gives $V_0=b_0\,r_{R,0}\,r_{C,0}$, and finally $\nu_1=K_0V_0=b_0\,r_{R,0}\,r_{C,0}\,K_0$. \end{proof}
